\chapter{Technical Appendix}
\label{chap:appendices}

\section{Reproducibility artifacts}
\label{app:reproducibility}

The main report contains the database schema, feature roles, evaluation metrics,
candidate portfolios, nested-evaluation procedure, uncertainty methods, and
model-selection policy. This appendix records only the technical identifiers
required to reproduce the reported artifacts.

The full source code, configuration files, report source, figure-generation
scripts, and persisted project structure are available in the public GitHub
repository:

\begin{center}
\url{https://github.com/manef42/NYU}
\end{center}

The repository contains separate directories for the ignition-classification
and FSR-regression pipelines, together with the \texttt{Report} directory that
contains the \LaTeX{} source files and figures used in this document.

The following identifiers refer to the modelling datasets used to generate the
persisted evaluation artifacts:

\begin{center}
\scriptsize
\texttt{FSR: 6e09a99e4ea69148e69151bd46149be4140d181b2b8294106927ea9499002fc8}\\
\texttt{Ignition: d80da0648ca3a7d77d31ae90ebea53b4e81b2ab28c8ec666ce6cd150ee7544f6}
\end{center}

For both tasks, the saved artifacts include:

\begin{itemize}
  \item cleaned modelling datasets and their SHA-256 fingerprints;
  \item frozen outer and inner split assignments;
  \item candidate configuration files and selected hyperparameters;
  \item outer-fold predictions and fold-level metrics;
  \item paper-level metrics where applicable;
  \item integrity checks, failure records, and selection outputs;
  \item refitted inference artifacts, model cards, and feature manifests.
\end{itemize}

Random seeds begin at 42. Evaluation artifacts and refitted artifacts have
different roles. Outer-fold predictions are the comparison evidence because
they are generated on untouched test partitions. Refit artifacts are trained
on all labelled rows and are intended for inference only; their training fit is
not used as evidence of generalisation.

\newpage
\section{FSR bootstrap augmentation example}

\Cref{fig:app-fsr-bootstrap-distribution} illustrates the bootstrap-augmentation
procedure for one FSR outer-training partition. The \texttt{rd} condition uses
the 2,081 real training observations. The \texttt{rd\_bt} condition keeps those
same observations and appends 1,000 rows sampled with replacement from them,
giving 3,081 training rows.

The augmented distribution preserves the original low-FSR-dominated shape
because bootstrap augmentation resamples existing observations rather than
creating new physical conditions. Some observations occur more often and others
less often in the augmented training set, but no new material, rig, flow,
gravity, or high-FSR regime is introduced.

\begin{figure}[htbp]
  \centering
  \includegraphics[width=0.82\textwidth]{fsr_bootstrapping.png}
  \caption{Example of FSR bootstrap augmentation for one outer-training
  partition. Blue bars show the real training data
  (\(n=2{,}081\)). Orange bars show the same real data plus 1,000
  bootstrap-resampled rows (\(n=3{,}081\)). Bootstrap augmentation changes
  the frequency of observed FSR values but does not create new physical
  support or enrich the sparse high-FSR tail.}
  \label{fig:app-fsr-bootstrap-distribution}
\end{figure}